\begin{frame}{Week 5.3의 중요도 샘플링과 무엇이 다른가}
  \begin{columns}[T]
    \begin{column}{0.5\textwidth}
      \kb{Week 5.3 (도구로 쓰일 때)}
      \begin{itemize}
        \item off-policy: target 정책의 데이터를 behavior 정책의
          데이터로 근사.
        \item ``무게가 크게 부등한 샘플들이 \textbf{편향 없이
          평균은 맞히되}, 작은 데이터에서는 한 샘플에 지배당한다'' —
          평균만 맞히면 됐다.
      \end{itemize}
    \end{column}
    \begin{column}{0.5\textwidth}
      \kb{여기서 (목적함수로 끌어다 쓸 때)}
      \begin{itemize}
        \item $L^{\mathrm{IS}}=\mathbb{E}[r_tA_t]$ 를 최대화
          \textbf{대상}으로.
        \item 같은 성질이 \textbf{학습의 불안정성}으로 직접 연결 —
          민감한 것은 ``평균''이 아니라 \textbf{매 스텝의
          그래디언트 방향}.
      \end{itemize}
    \end{column}
  \end{columns}
  \vskip0.8em
  \begin{block}{이 절의 통찰}
    ``Week 5.3의 도구가 여기서는 \textbf{안전문제}가 된다.''
  \end{block}
\end{frame}
